\clearpage

\section{绪论}

\subsection{研究背景与意义}

学位论文是研究生科研工作的最终凝结，其排版规范不仅关乎论文的视觉呈现，也影响评审人对作者学术态度的初步判断。深圳大学《研究生学位论文基本要求及格式规范》对纸张、页面、字体、章节、图表与参考文献等环节均做出了细致约束\footnote{本模板对照的版本为研究生院于2024年1月发布的最新规范。}，作者若以人工方式逐项核对，需要付出可观的时间成本。

由于学校官方仅以Word文档形式发布模板，作者在使用Word撰写论文时，常面临大量排版问题。LaTeX凭借稳定的版式控制能力，长期以来都是学位论文写作的更优选择；但LaTeX的学习曲线陡峭，从零搭建一份合规的论文工程，对绝大多数研究生而言仍是不小的负担。本模板正是在这一背景下产生\footnote{本模板最初源自\url{https://github.com/Ftine/SZU_Latex_Master_Template}，原作者已停止维护，本仓库接手并持续迭代。}：以一份对齐学校最新规范、开箱即用的工程，让作者将精力集中于研究本身。

\subsection{本模板的设计理念}

\subsubsection{自文档化}

本模板采用“自文档化”的设计：正文各章节本身即为模板的使用说明。读者克隆仓库后直接编译生成的初始PDF，将依次呈现摘要的写作要点、绪论的组织方式、各级标题的层级关系、公式与图表的编号规则、参考文献的引用方式，以及致谢与研究成果的标准格式。在阅读PDF的过程中，作者即可同步理解每一处排版背后所对应的学校规范条款；着手写作时，只需将示范文字逐节替换为自己的研究内容，模板的格式约束将自动延续。

这一理念的好处在于：模板的“文档”与“实现”天然同步。当学校规范更新、模板代码相应调整时，正文示范也会随之改写，作者所看到的始终是当前实现下最准确的使用方式，不存在文档与代码脱节的风险。

\subsubsection{兼容论文归档流程}

学位论文从写作到最终归档需经历预答辩、盲审、答辩等多个环节。在写作与盲审阶段，正文无需附带学术评语、答辩委员会决议书等评审材料；但论文最终归档时，这些材料须作为附件随论文一并提交。

本模板对这一流程做了预设：主控文件\texttt{master\_pang.tex}的末尾保留了若干被注释的代码片段，分别用于在目录中追加“学术评语”与“答辩委员会决议书”条目，以及通过\verb|\includepdf| 将这两份PDF文件插入论文末尾。作者在写作期间无须关心这些代码；待归档环节临近，将相应行的注释符号去除并放置好对应的PDF文件，即可一次性完成归档版本的生成，无需对模板做侵入性修改。

\subsubsection{专硕与学硕通用}

学校对专业硕士与学术硕士学位论文的格式要求一致，差异仅体现在封面信息上。本模板将封面拆分为独立的PDF文件\texttt{chapter/masterPage.pdf}，由主控文件通过\verb|\includepdf|引入，正文结构与各项排版约束完全独立于封面。作者只需依据自身学位类型替换该PDF文件，即可在专硕与学硕之间无缝切换，其余源文件无须任何改动。事实上，学校对博士学位论文的格式要求与硕士一致，本模板后续也计划将适用范围扩展至博士论文，进一步提升通用性。

\subsection{本模板的组织结构}

本模板的源文件围绕主控文件\texttt{master\_pang.tex}组织，宏包加载、页面与字体设置、各类编号约定均集中于该文件。正文按章节拆分至\texttt{chapter/}目录下的独立\texttt{tex}文件，由主控文件通过\verb|\input|逐一引入。这样的结构既便于作者按章并行写作，也使版本控制工具能够更清晰地追踪每一章的演进。

除正文外，模板还包含若干辅助目录：\texttt{figure/}用于按章节存放插图源文件，\texttt{tables/}用于将复杂表格的源码独立维护，\texttt{ref\_paper.bib}集中保存参考文献条目。后续各章将依次示范这些元素的具体用法。
